
\documentclass{article}
\usepackage{colortbl}
\usepackage{makecell}
\usepackage{multirow}
\usepackage{supertabular}

\begin{document}

\newcounter{utterance}

\centering \large Interaction Transcript for game `matchit\_ascii', experiment `different\_grid', episode 0 with Apertus{-}v1.5{-}8B.
\vspace{24pt}

{ \footnotesize  \setcounter{utterance}{1}
\setlength{\tabcolsep}{0pt}
\begin{supertabular}{c@{$\;$}|p{.15\linewidth}@{}p{.15\linewidth}p{.15\linewidth}p{.15\linewidth}p{.15\linewidth}p{.15\linewidth}}
    \# & $\;$A & \multicolumn{4}{c}{Game Master} & $\;\:$B\\
    \hline

    \theutterance \stepcounter{utterance}  
    & & \multicolumn{4}{p{0.6\linewidth}}{
        \cellcolor[rgb]{0.9,0.9,0.9}{
            \makecell[{{p{\linewidth}}}]{
                \texttt{\tiny{[P1$\langle$GM]}}
                \texttt{Συμμετέχετε σε ένα συνεργατικό παιχνίδι μαντεψιάς.} \\
\texttt{Ο στόχος είναι να βρείτε αν αυτό το πλέγμα και μια άλλη εικόνα που μόνο εγώ μπορώ να δω, είναι τα ίδια.} \\
\texttt{Παρακαλώ περιγράψτε πρώτα το πλέγμα σας. Στη συνέχεια, θα δώσω τη δική μου περιγραφή και μπορούμε να κάνουμε ο ένας στον άλλο ερωτήσεις σχετικά με τα πλέγματα για να καταλάβουμε αν είναι τα ίδια. Ξεκινήστε τώρα τη σύντομη περιγραφή του πλέγματός σας με "ΠΕΡΙΓΡΑΦΗ:" ακολουθούμενη από την περιγραφή. Μην προσθέσετε τίποτα άλλο. Αυτό είναι το πλέγμα σας:} \\
\texttt{X ▢ ▢ ▢ X} \\
\texttt{X ▢ ▢ ▢ X} \\
\texttt{X ▢ X ▢ X} \\
\texttt{X X ▢ X X} \\
\texttt{X ▢ ▢ ▢ X} \\
            }
        }
    }
    & \\ \\

    \theutterance \stepcounter{utterance}  
    & \multicolumn{4}{p{0.6\linewidth}}{
        \cellcolor[rgb]{0.9,0.9,0.9}{
            \makecell[{{p{\linewidth}}}]{
                \texttt{\tiny{[P1$\rangle$GM]}}
                \texttt{ΠΕΡΙΓΡΑΦΗ: Ένα ορθογώνιο πλέγμα 4 σειρών και 5 στηλών με 'X' και κενά (▢) σε διάφορες θέσεις.} \\
            }
        }
    }
    & & \\ \\

    \theutterance \stepcounter{utterance}  
    & & & \multicolumn{2}{p{0.3\linewidth}}{
        \cellcolor[rgb]{0.9,0.9,0.9}{
            \makecell[{{p{\linewidth}}}]{
                \texttt{\tiny{[GM$|$GM]}}
                \texttt{continue} \\
            }
        }
    }
    & & \\ \\

    \theutterance \stepcounter{utterance}  
    & & \multicolumn{4}{p{0.6\linewidth}}{
        \cellcolor[rgb]{0.9,0.9,0.9}{
            \makecell[{{p{\linewidth}}}]{
                \texttt{\tiny{[GM$\rangle$P2]}}
                \texttt{Συμμετέχετε σε ένα συνεργατικό παιχνίδι μαντεψιάς.} \\
\texttt{Ο στόχος είναι να βρείτε αν αυτό το πλέγμα και μια άλλη εικόνα που μόνο εγώ μπορώ να δω, είναι τα ίδια.} \\
\texttt{Παρακαλώ περιγράψτε πρώτα το πλέγμα σας. Στη συνέχεια, θα δώσω τη δική μου περιγραφή και μπορούμε να κάνουμε ο ένας στον άλλο ερωτήσεις σχετικά με τα πλέγματα για να καταλάβουμε αν είναι τα ίδια. Ξεκινήστε τώρα τη σύντομη περιγραφή του πλέγματός σας με "ΠΕΡΙΓΡΑΦΗ:" ακολουθούμενη από την περιγραφή. Μην προσθέσετε τίποτα άλλο. Αυτό είναι το πλέγμα σας:} \\
\texttt{X X X X X} \\
\texttt{X ▢ ▢ ▢ X} \\
\texttt{X ▢ ▢ ▢ X} \\
\texttt{X ▢ ▢ ▢ X} \\
\texttt{X X X X X} \\
            }
        }
    }
    & \\ \\

    \theutterance \stepcounter{utterance}  
    & & & \multicolumn{4}{p{0.6\linewidth}}{
        \cellcolor[rgb]{0.9,0.9,0.9}{
            \makecell[{{p{\linewidth}}}]{
                \texttt{\tiny{[GM$\langle$P2]}}
                \texttt{ΠΕΡΙΓΡΑΦΗ: Το πλέγμα αποτελείται από 5 σειρές και 5 στήλες. Η πρώτη και τελευταία σειρά περιέχουν μόνο X σε κάθε θέση. Η δεύτερη, τρίτη και τέταρτη σειρά έχουν X στις δύο άκρες και ▢ (κύκλο) στο κέντρο.} \\
            }
        }
    }
     \\ \\

    \theutterance \stepcounter{utterance}  
    & & & \multicolumn{2}{p{0.3\linewidth}}{
        \cellcolor[rgb]{0.9,0.9,0.9}{
            \makecell[{{p{\linewidth}}}]{
                \texttt{\tiny{[GM$|$GM]}}
                \texttt{continue} \\
            }
        }
    }
    & & \\ \\

    \theutterance \stepcounter{utterance}  
    & & \multicolumn{4}{p{0.6\linewidth}}{
        \cellcolor[rgb]{0.9,0.9,0.9}{
            \makecell[{{p{\linewidth}}}]{
                \texttt{\tiny{[GM$\rangle$P2]}}
                \texttt{Αυτή είναι η ΠΕΡΙΓΡΑΦΗ μουΠΕΡΙΓΡΑΦΗ: Ένα ορθογώνιο πλέγμα 4 σειρών και 5 στηλών με 'X' και κενά (▢) σε διάφορες θέσεις.} \\
\texttt{Τώρα κάντε μια ερώτηση για να μάθετε νέες πτυχές του πλέγματός μου που μπορεί να διαφέρουν από το δικό σας πλέγμα. Ξεκινήστε με "ΕΡΩΤΗΣΗ:" και μην προσθέσετε τίποτα άλλο.} \\
            }
        }
    }
    & \\ \\

    \theutterance \stepcounter{utterance}  
    & & & \multicolumn{4}{p{0.6\linewidth}}{
        \cellcolor[rgb]{0.9,0.9,0.9}{
            \makecell[{{p{\linewidth}}}]{
                \texttt{\tiny{[GM$\langle$P2]}}
                \texttt{ΕΡΩΤΗΣΗ: Ποια είναι η ακριβής θέση των κενών (▢) στο πλέγμα σας;} \\
            }
        }
    }
     \\ \\

    \theutterance \stepcounter{utterance}  
    & & & \multicolumn{2}{p{0.3\linewidth}}{
        \cellcolor[rgb]{0.9,0.9,0.9}{
            \makecell[{{p{\linewidth}}}]{
                \texttt{\tiny{[GM$|$GM]}}
                \texttt{continue} \\
            }
        }
    }
    & & \\ \\

    \theutterance \stepcounter{utterance}  
    & & \multicolumn{4}{p{0.6\linewidth}}{
        \cellcolor[rgb]{0.9,0.9,0.9}{
            \makecell[{{p{\linewidth}}}]{
                \texttt{\tiny{[P1$\langle$GM]}}
                \texttt{Αυτή είναι η ΠΕΡΙΓΡΑΦΗ μουΠΕΡΙΓΡΑΦΗ: Το πλέγμα αποτελείται από 5 σειρές και 5 στήλες. Η πρώτη και τελευταία σειρά περιέχουν μόνο X σε κάθε θέση. Η δεύτερη, τρίτη και τέταρτη σειρά έχουν X στις δύο άκρες και ▢ (κύκλο) στο κέντρο.} \\
\texttt{ΕΡΩΤΗΣΗ: Ποια είναι η ακριβής θέση των κενών (▢) στο πλέγμα σας;Ξεκινήστε την απάντησή σας με "ΑΠΑΝΤΗΣΗ:" και μην προσθέσετε τίποτα άλλο.} \\
            }
        }
    }
    & \\ \\

    \theutterance \stepcounter{utterance}  
    & \multicolumn{4}{p{0.6\linewidth}}{
        \cellcolor[rgb]{0.9,0.9,0.9}{
            \makecell[{{p{\linewidth}}}]{
                \texttt{\tiny{[P1$\rangle$GM]}}
                \texttt{ΑΠΑΝΤΗΣΗ: Η θέση των κενών (▢) είναι στη δεύτερη, τρίτη και τέταρτη σειρά, συγκεκριμένα στην τρίτη θέση κάθε μιας από αυτές τις σειρές.} \\
            }
        }
    }
    & & \\ \\

    \theutterance \stepcounter{utterance}  
    & & & \multicolumn{2}{p{0.3\linewidth}}{
        \cellcolor[rgb]{0.9,0.9,0.9}{
            \makecell[{{p{\linewidth}}}]{
                \texttt{\tiny{[GM$|$GM]}}
                \texttt{continue} \\
            }
        }
    }
    & & \\ \\

    \theutterance \stepcounter{utterance}  
    & & \multicolumn{4}{p{0.6\linewidth}}{
        \cellcolor[rgb]{0.9,0.9,0.9}{
            \makecell[{{p{\linewidth}}}]{
                \texttt{\tiny{[P1$\langle$GM]}}
                \texttt{Τώρα κάντε μια ερώτηση για να μάθετε νέες πτυχές του πλέγματός μου που μπορεί να διαφέρουν από το δικό σας πλέγμα. Ξεκινήστε με "ΕΡΩΤΗΣΗ:" και μην προσθέσετε τίποτα άλλο.} \\
            }
        }
    }
    & \\ \\

    \theutterance \stepcounter{utterance}  
    & \multicolumn{4}{p{0.6\linewidth}}{
        \cellcolor[rgb]{0.9,0.9,0.9}{
            \makecell[{{p{\linewidth}}}]{
                \texttt{\tiny{[P1$\rangle$GM]}}
                \texttt{ΕΡΩΤΗΣΗ: Υπάρχουν X σε οποιαδήποτε θέση της τελευταίας (πέμπτης) σειράς εκτός από τις άκρες;} \\
            }
        }
    }
    & & \\ \\

    \theutterance \stepcounter{utterance}  
    & & & \multicolumn{2}{p{0.3\linewidth}}{
        \cellcolor[rgb]{0.9,0.9,0.9}{
            \makecell[{{p{\linewidth}}}]{
                \texttt{\tiny{[GM$|$GM]}}
                \texttt{continue} \\
            }
        }
    }
    & & \\ \\

    \theutterance \stepcounter{utterance}  
    & & \multicolumn{4}{p{0.6\linewidth}}{
        \cellcolor[rgb]{0.9,0.9,0.9}{
            \makecell[{{p{\linewidth}}}]{
                \texttt{\tiny{[GM$\rangle$P2]}}
                \texttt{ΑΠΑΝΤΗΣΗ: Η θέση των κενών (▢) είναι στη δεύτερη, τρίτη και τέταρτη σειρά, συγκεκριμένα στην τρίτη θέση κάθε μιας από αυτές τις σειρές.} \\
\texttt{ΕΡΩΤΗΣΗ: Υπάρχουν X σε οποιαδήποτε θέση της τελευταίας (πέμπτης) σειράς εκτός από τις άκρες;Ξεκινήστε την απάντησή σας με "ΑΠΑΝΤΗΣΗ:" και μην προσθέσετε τίποτα άλλο.} \\
            }
        }
    }
    & \\ \\

    \theutterance \stepcounter{utterance}  
    & & & \multicolumn{4}{p{0.6\linewidth}}{
        \cellcolor[rgb]{0.9,0.9,0.9}{
            \makecell[{{p{\linewidth}}}]{
                \texttt{\tiny{[GM$\langle$P2]}}
                \texttt{ΑΠΑΝΤΗΣΗ: Όχι, η πέμπτη σειρά περιέχει μόνο X στις δύο άκρες και δεν υπάρχουν άλλα στοιχεία στο κέντρο.} \\
            }
        }
    }
     \\ \\

    \theutterance \stepcounter{utterance}  
    & & & \multicolumn{2}{p{0.3\linewidth}}{
        \cellcolor[rgb]{0.9,0.9,0.9}{
            \makecell[{{p{\linewidth}}}]{
                \texttt{\tiny{[GM$|$GM]}}
                \texttt{continue} \\
            }
        }
    }
    & & \\ \\

    \theutterance \stepcounter{utterance}  
    & & \multicolumn{4}{p{0.6\linewidth}}{
        \cellcolor[rgb]{0.9,0.9,0.9}{
            \makecell[{{p{\linewidth}}}]{
                \texttt{\tiny{[GM$\rangle$P2]}}
                \texttt{Τώρα κάντε μια ερώτηση για να μάθετε νέες πτυχές του πλέγματός μου που μπορεί να διαφέρουν από το δικό σας πλέγμα. Ξεκινήστε με "ΕΡΩΤΗΣΗ:" και μην προσθέσετε τίποτα άλλο.} \\
            }
        }
    }
    & \\ \\

    \theutterance \stepcounter{utterance}  
    & & & \multicolumn{4}{p{0.6\linewidth}}{
        \cellcolor[rgb]{0.9,0.9,0.9}{
            \makecell[{{p{\linewidth}}}]{
                \texttt{\tiny{[GM$\langle$P2]}}
                \texttt{ΕΡΩΤΗΣΗ: Υπάρχει οποιαδήποτε άλλη μη X θέση στο πλέγμα εκτός από αυτά που έχω ήδη περιγράψει;} \\
            }
        }
    }
     \\ \\

    \theutterance \stepcounter{utterance}  
    & & & \multicolumn{2}{p{0.3\linewidth}}{
        \cellcolor[rgb]{0.9,0.9,0.9}{
            \makecell[{{p{\linewidth}}}]{
                \texttt{\tiny{[GM$|$GM]}}
                \texttt{continue} \\
            }
        }
    }
    & & \\ \\

    \theutterance \stepcounter{utterance}  
    & & \multicolumn{4}{p{0.6\linewidth}}{
        \cellcolor[rgb]{0.9,0.9,0.9}{
            \makecell[{{p{\linewidth}}}]{
                \texttt{\tiny{[P1$\langle$GM]}}
                \texttt{ΑΠΑΝΤΗΣΗ: Όχι, η πέμπτη σειρά περιέχει μόνο X στις δύο άκρες και δεν υπάρχουν άλλα στοιχεία στο κέντρο.} \\
\texttt{ΕΡΩΤΗΣΗ: Υπάρχει οποιαδήποτε άλλη μη X θέση στο πλέγμα εκτός από αυτά που έχω ήδη περιγράψει;Ξεκινήστε την απάντησή σας με "ΑΠΑΝΤΗΣΗ:" και μην προσθέσετε τίποτα άλλο.} \\
            }
        }
    }
    & \\ \\

    \theutterance \stepcounter{utterance}  
    & \multicolumn{4}{p{0.6\linewidth}}{
        \cellcolor[rgb]{0.9,0.9,0.9}{
            \makecell[{{p{\linewidth}}}]{
                \texttt{\tiny{[P1$\rangle$GM]}}
                \texttt{ΑΠΑΝΤΗΣΗ: Όχι, δεν υπάρχουν άλλες μη X θέσεις στο πλέγμα εκτός από αυτές που έχουν ήδη περιγραφεί.} \\
            }
        }
    }
    & & \\ \\

    \theutterance \stepcounter{utterance}  
    & & & \multicolumn{2}{p{0.3\linewidth}}{
        \cellcolor[rgb]{0.9,0.9,0.9}{
            \makecell[{{p{\linewidth}}}]{
                \texttt{\tiny{[GM$|$GM]}}
                \texttt{continue} \\
            }
        }
    }
    & & \\ \\

    \theutterance \stepcounter{utterance}  
    & & \multicolumn{4}{p{0.6\linewidth}}{
        \cellcolor[rgb]{0.9,0.9,0.9}{
            \makecell[{{p{\linewidth}}}]{
                \texttt{\tiny{[P1$\langle$GM]}}
                \texttt{Τώρα κάντε μια ερώτηση για να μάθετε νέες πτυχές του πλέγματός μου που μπορεί να διαφέρουν από το δικό σας πλέγμα. Ξεκινήστε με "ΕΡΩΤΗΣΗ:" και μην προσθέσετε τίποτα άλλο.} \\
            }
        }
    }
    & \\ \\

    \theutterance \stepcounter{utterance}  
    & \multicolumn{4}{p{0.6\linewidth}}{
        \cellcolor[rgb]{0.9,0.9,0.9}{
            \makecell[{{p{\linewidth}}}]{
                \texttt{\tiny{[P1$\rangle$GM]}}
                \texttt{ΕΡΩΤΗΣΗ: Το πλέγμα περιέχει μόνο X και ▢, ή υπάρχουν και άλλα σύμβολα ή διαφορετικά στοιχεία;} \\
            }
        }
    }
    & & \\ \\

    \theutterance \stepcounter{utterance}  
    & & & \multicolumn{2}{p{0.3\linewidth}}{
        \cellcolor[rgb]{0.9,0.9,0.9}{
            \makecell[{{p{\linewidth}}}]{
                \texttt{\tiny{[GM$|$GM]}}
                \texttt{continue} \\
            }
        }
    }
    & & \\ \\

    \theutterance \stepcounter{utterance}  
    & & \multicolumn{4}{p{0.6\linewidth}}{
        \cellcolor[rgb]{0.9,0.9,0.9}{
            \makecell[{{p{\linewidth}}}]{
                \texttt{\tiny{[GM$\rangle$P2]}}
                \texttt{ΑΠΑΝΤΗΣΗ: Όχι, δεν υπάρχουν άλλες μη X θέσεις στο πλέγμα εκτός από αυτές που έχουν ήδη περιγραφεί.} \\
\texttt{ΕΡΩΤΗΣΗ: Το πλέγμα περιέχει μόνο X και ▢, ή υπάρχουν και άλλα σύμβολα ή διαφορετικά στοιχεία;Ξεκινήστε την απάντησή σας με "ΑΠΑΝΤΗΣΗ:" και μην προσθέσετε τίποτα άλλο.} \\
            }
        }
    }
    & \\ \\

    \theutterance \stepcounter{utterance}  
    & & & \multicolumn{4}{p{0.6\linewidth}}{
        \cellcolor[rgb]{0.9,0.9,0.9}{
            \makecell[{{p{\linewidth}}}]{
                \texttt{\tiny{[GM$\langle$P2]}}
                \texttt{ΑΠΑΝΤΗΣΗ: Το πλέγμα περιέχει μόνο X και ▢, δεν υπάρχουν άλλα σύμβολα ή διαφορετικά στοιχεία.} \\
            }
        }
    }
     \\ \\

    \theutterance \stepcounter{utterance}  
    & & & \multicolumn{2}{p{0.3\linewidth}}{
        \cellcolor[rgb]{0.9,0.9,0.9}{
            \makecell[{{p{\linewidth}}}]{
                \texttt{\tiny{[GM$|$GM]}}
                \texttt{continue} \\
            }
        }
    }
    & & \\ \\

    \theutterance \stepcounter{utterance}  
    & & \multicolumn{4}{p{0.6\linewidth}}{
        \cellcolor[rgb]{0.9,0.9,0.9}{
            \makecell[{{p{\linewidth}}}]{
                \texttt{\tiny{[GM$\rangle$P2]}}
                \texttt{Τώρα κάντε μια ερώτηση για να μάθετε νέες πτυχές του πλέγματός μου που μπορεί να διαφέρουν από το δικό σας πλέγμα. Ξεκινήστε με "ΕΡΩΤΗΣΗ:" και μην προσθέσετε τίποτα άλλο.} \\
            }
        }
    }
    & \\ \\

    \theutterance \stepcounter{utterance}  
    & & & \multicolumn{4}{p{0.6\linewidth}}{
        \cellcolor[rgb]{0.9,0.9,0.9}{
            \makecell[{{p{\linewidth}}}]{
                \texttt{\tiny{[GM$\langle$P2]}}
                \texttt{ΕΡΩΤΗΣΗ: Το πλέγμα έχει συμμετρική σχέση μεταξύ των σειρών ή είναι εντελώς ασύμμετρο;} \\
            }
        }
    }
     \\ \\

    \theutterance \stepcounter{utterance}  
    & & & \multicolumn{2}{p{0.3\linewidth}}{
        \cellcolor[rgb]{0.9,0.9,0.9}{
            \makecell[{{p{\linewidth}}}]{
                \texttt{\tiny{[GM$|$GM]}}
                \texttt{continue} \\
            }
        }
    }
    & & \\ \\

    \theutterance \stepcounter{utterance}  
    & & \multicolumn{4}{p{0.6\linewidth}}{
        \cellcolor[rgb]{0.9,0.9,0.9}{
            \makecell[{{p{\linewidth}}}]{
                \texttt{\tiny{[P1$\langle$GM]}}
                \texttt{ΑΠΑΝΤΗΣΗ: Το πλέγμα περιέχει μόνο X και ▢, δεν υπάρχουν άλλα σύμβολα ή διαφορετικά στοιχεία.} \\
\texttt{ΕΡΩΤΗΣΗ: Το πλέγμα έχει συμμετρική σχέση μεταξύ των σειρών ή είναι εντελώς ασύμμετρο;Ξεκινήστε την απάντησή σας με "ΑΠΑΝΤΗΣΗ:" και μην προσθέσετε τίποτα άλλο.} \\
            }
        }
    }
    & \\ \\

    \theutterance \stepcounter{utterance}  
    & \multicolumn{4}{p{0.6\linewidth}}{
        \cellcolor[rgb]{0.9,0.9,0.9}{
            \makecell[{{p{\linewidth}}}]{
                \texttt{\tiny{[P1$\rangle$GM]}}
                \texttt{ΑΠΑΝΤΗΣΗ: Το πλέγμα δεν είναι συμμετρικό μεταξύ των σειρών.} \\
            }
        }
    }
    & & \\ \\

    \theutterance \stepcounter{utterance}  
    & & & \multicolumn{2}{p{0.3\linewidth}}{
        \cellcolor[rgb]{0.9,0.9,0.9}{
            \makecell[{{p{\linewidth}}}]{
                \texttt{\tiny{[GM$|$GM]}}
                \texttt{continue} \\
            }
        }
    }
    & & \\ \\

    \theutterance \stepcounter{utterance}  
    & & \multicolumn{4}{p{0.6\linewidth}}{
        \cellcolor[rgb]{0.9,0.9,0.9}{
            \makecell[{{p{\linewidth}}}]{
                \texttt{\tiny{[P1$\langle$GM]}}
                \texttt{Τώρα κάντε μια ερώτηση για να μάθετε νέες πτυχές του πλέγματός μου που μπορεί να διαφέρουν από το δικό σας πλέγμα. Ξεκινήστε με "ΕΡΩΤΗΣΗ:" και μην προσθέσετε τίποτα άλλο.} \\
            }
        }
    }
    & \\ \\

    \theutterance \stepcounter{utterance}  
    & \multicolumn{4}{p{0.6\linewidth}}{
        \cellcolor[rgb]{0.9,0.9,0.9}{
            \makecell[{{p{\linewidth}}}]{
                \texttt{\tiny{[P1$\rangle$GM]}}
                \texttt{ΕΡΩΤΗΣΗ: Υπάρχουν διαγώνια σχέσεις ή μοτίβα που συνδέουν συγκεκριμένες θέσεις μεταξύ των σειρών ή των στηλών;} \\
            }
        }
    }
    & & \\ \\

    \theutterance \stepcounter{utterance}  
    & & & \multicolumn{2}{p{0.3\linewidth}}{
        \cellcolor[rgb]{0.9,0.9,0.9}{
            \makecell[{{p{\linewidth}}}]{
                \texttt{\tiny{[GM$|$GM]}}
                \texttt{continue} \\
            }
        }
    }
    & & \\ \\

    \theutterance \stepcounter{utterance}  
    & & \multicolumn{4}{p{0.6\linewidth}}{
        \cellcolor[rgb]{0.9,0.9,0.9}{
            \makecell[{{p{\linewidth}}}]{
                \texttt{\tiny{[GM$\rangle$P2]}}
                \texttt{ΑΠΑΝΤΗΣΗ: Το πλέγμα δεν είναι συμμετρικό μεταξύ των σειρών.} \\
\texttt{ΕΡΩΤΗΣΗ: Υπάρχουν διαγώνια σχέσεις ή μοτίβα που συνδέουν συγκεκριμένες θέσεις μεταξύ των σειρών ή των στηλών;Ξεκινήστε την απάντησή σας με "ΑΠΑΝΤΗΣΗ:" και μην προσθέσετε τίποτα άλλο.} \\
            }
        }
    }
    & \\ \\

    \theutterance \stepcounter{utterance}  
    & & & \multicolumn{4}{p{0.6\linewidth}}{
        \cellcolor[rgb]{0.9,0.9,0.9}{
            \makecell[{{p{\linewidth}}}]{
                \texttt{\tiny{[GM$\langle$P2]}}
                \texttt{ΑΠΑΝΤΗΣΗ: Δεν υπάρχουν διαγώνια σχέσεις ή μοτίβα που συνδέουν συγκεκριμένες θέσεις μεταξύ των σειρών ή των στηλών.} \\
            }
        }
    }
     \\ \\

    \theutterance \stepcounter{utterance}  
    & & & \multicolumn{2}{p{0.3\linewidth}}{
        \cellcolor[rgb]{0.9,0.9,0.9}{
            \makecell[{{p{\linewidth}}}]{
                \texttt{\tiny{[GM$|$GM]}}
                \texttt{continue} \\
            }
        }
    }
    & & \\ \\

    \theutterance \stepcounter{utterance}  
    & & \multicolumn{4}{p{0.6\linewidth}}{
        \cellcolor[rgb]{0.9,0.9,0.9}{
            \makecell[{{p{\linewidth}}}]{
                \texttt{\tiny{[GM$\rangle$P2]}}
                \texttt{Τώρα πάρε μια απόφαση. Τι πιστεύεις: είναι το δικό σου πλέγμα και το άλλο πλέγμα που περιγράφεται το ίδιο; Γράψε "ΑΠΟΦΑΣΗ: ίδιο πλέγμα" αν πιστεύεις ότι είναι το ίδιο πλέγμα ή "ΑΠΟΦΑΣΗ: διαφορετικά πλέγματα" αν πιστεύεις ότι είναι διαφορετικά πλέγματα. Μην προσθέσεις τίποτα άλλο.} \\
            }
        }
    }
    & \\ \\

    \theutterance \stepcounter{utterance}  
    & & & \multicolumn{4}{p{0.6\linewidth}}{
        \cellcolor[rgb]{0.9,0.9,0.9}{
            \makecell[{{p{\linewidth}}}]{
                \texttt{\tiny{[GM$\langle$P2]}}
                \texttt{ΑΠΟΦΑΣΗ: ίδιο πλέγμα} \\
            }
        }
    }
     \\ \\

    \theutterance \stepcounter{utterance}  
    & & & \multicolumn{2}{p{0.3\linewidth}}{
        \cellcolor[rgb]{0.9,0.9,0.9}{
            \makecell[{{p{\linewidth}}}]{
                \texttt{\tiny{[GM$|$GM]}}
                \texttt{continue} \\
            }
        }
    }
    & & \\ \\

    \theutterance \stepcounter{utterance}  
    & & & \multicolumn{2}{p{0.3\linewidth}}{
        \cellcolor[rgb]{0.9,0.9,0.9}{
            \makecell[{{p{\linewidth}}}]{
                \texttt{\tiny{[GM$|$GM]}}
                \texttt{loss} \\
            }
        }
    }
    & & \\ \\

    \theutterance \stepcounter{utterance}  
    & & \multicolumn{4}{p{0.6\linewidth}}{
        \cellcolor[rgb]{0.9,0.9,0.9}{
            \makecell[{{p{\linewidth}}}]{
                \texttt{\tiny{[P1$\langle$GM]}}
                \texttt{ΑΠΑΝΤΗΣΗ: Δεν υπάρχουν διαγώνια σχέσεις ή μοτίβα που συνδέουν συγκεκριμένες θέσεις μεταξύ των σειρών ή των στηλών.} \\
\texttt{Τώρα πάρε μια απόφαση. Τι πιστεύεις: είναι το δικό σου πλέγμα και το άλλο πλέγμα που περιγράφεται το ίδιο; Γράψε "ΑΠΟΦΑΣΗ: ίδιο πλέγμα" αν πιστεύεις ότι είναι το ίδιο πλέγμα ή "ΑΠΟΦΑΣΗ: διαφορετικά πλέγματα" αν πιστεύεις ότι είναι διαφορετικά πλέγματα. Μην προσθέσεις τίποτα άλλο.} \\
            }
        }
    }
    & \\ \\

    \theutterance \stepcounter{utterance}  
    & \multicolumn{4}{p{0.6\linewidth}}{
        \cellcolor[rgb]{0.9,0.9,0.9}{
            \makecell[{{p{\linewidth}}}]{
                \texttt{\tiny{[P1$\rangle$GM]}}
                \texttt{ΑΠΟΦΑΣΗ: ίδια πλέγματα} \\
            }
        }
    }
    & & \\ \\

    \theutterance \stepcounter{utterance}  
    & & & \multicolumn{2}{p{0.3\linewidth}}{
        \cellcolor[rgb]{0.9,0.9,0.9}{
            \makecell[{{p{\linewidth}}}]{
                \texttt{\tiny{[GM$|$GM]}}
                \texttt{continue} \\
            }
        }
    }
    & & \\ \\

    \theutterance \stepcounter{utterance}  
    & & & \multicolumn{2}{p{0.3\linewidth}}{
        \cellcolor[rgb]{0.9,0.9,0.9}{
            \makecell[{{p{\linewidth}}}]{
                \texttt{\tiny{[GM$|$GM]}}
                \texttt{abort, wrong message content.} \\
            }
        }
    }
    & & \\ \\

\end{supertabular}
}

\end{document}
